\documentclass[../../main.tex]{subfiles}
\usepackage{amsmath}	% required for `\align' (yatex added)
\begin{document}

{\bf [解]}
$F(x) = \int_0^x f(t) dt$ とおくと，具体的な$f(x)$の表式から
\begin{align}
F(x) = \dfrac{1}{4}x^4 + \dfrac{p}{3}x^3 + \dfrac{q}{2}x^2 + rx \label{eq:1}
\end{align}
である．与えられた不等式は$F(x)$を用いて
\begin{align}
   2F(x) \le F(x+y) + F(x-y) \label{eq:2}
\end{align}
と書ける．右辺に\cref{eq:1}を代入して
\begin{align*}
  \mathrm{RHS} 
  & = \frac{1}{4}\left[(x+y)^4 + (x-y)^4 \right] + \frac{p}{3}\left[(x+y)^3 + (x-y)^3 \right] 
  + \frac{q}{2}\left[(x+y)^2 + (x-y)^2 \right] + r\left[(x+y) + (x-y) \right] \\
  & = \frac{1}{4}\left[2x^4 + 12x^2 y^2 + 2y^4\right] + \frac{p}{3}\left[2x^3 + 6xy^2\right] + \frac{q}{2}\left[2x^2+2y^2\right] + 2rx
\end{align*}
だから左辺との差をとると
\begin{align}
 \mathrm{RHS}-\mathrm{LHS} 
 &= \frac{1}{4}\left[12x^2 y^2 + 2y^4\right] + \frac{p}{3}\left[6xy^2\right] + \frac{q}{2}\left[2y^2\right] \\
 &= y^2\left(3x^2 + \frac{1}{2}y^2 + 2px + q\right)
\end{align}
となる．これが非負になる条件を求めれば良い．

この式は$y=0$の時は$0$であり与えられた不等式\cref{eq:2}は成り立つ．$y\neq 0$の時は $y^2 > 0$だから両辺$y^2$で割ると
\begin{align}
 \left(3x^2 + 2px + q\right) + \frac{1}{2}y^2 \ge 0
\end{align}
が任意の$(x,y)$に対して成り立てば良い．そのためには$3x^2 + 2px + q = 0$ の判別式 $D$ として$D\le 0$であれば良い．
\begin{align}
  D \le 0 
  & \iff p^2 - 3q \le 0 \\
  &\iff p^2 \le 3q 
\end{align}
だから，求める条件は
\begin{align}
 p^2 \le 3q 
\end{align}
である．

{\bf [解説]}
最後の条件に$r$が出てこないのは自然で，$r$は$f(x)$の定数項だから，元の不等式で左辺で$[0,x]$まで積分したものの$2$倍と，
右辺で$[0,x+y]$,$[0,x-y]$まで積分したものの和は等しくなって必ず打ち消しあうので条件として現れない．

途中で出てきた式\eqref{eq:2}はまさに凸不等式であり，要するに$F(x)$が下に凸である条件を求めれば良い．
以下その方針にたった解答を示す．こちらの方が数学的な背景がわかってる感があって綺麗だろう．

$F(x)$が下に凸なための条件は2階微分が非負であること，すなわち$F''(x)\ge 0$である．実際計算すると
\begin{align}
 F''(x) = 3x^2+2px+q
\end{align}
だから$f'(x)=0$の判別式
\begin{align}
 D = p^2-3q
\end{align}
が非正であればよく，
\begin{align}
 p^2 \le 3p
\end{align}
が求める条件である．

\newpage

\end{document}
